\section{Introduction}
\label{submission:sec:intro}

Deep learning in the wild is increasingly transitioning from massive centralized cloud servers to decentralized, resource-constrained edge and IoT devices (e.g., smart cameras, mobile devices, autonomous drones). On these devices, deploying multiple independent, specialized foundation models—each parameterized with hundreds of millions of weights—is computationally, financially, and logistically prohibitive due to severe storage, memory bandwidth, and network transmission bottlenecks. For instance, storing five independent fine-tuned visual encoders from the CLIP family \cite{radford2021learning} can easily consume gigabytes of storage, rendering over-the-air updates or on-device expert switching completely unviable.

To mitigate this bottleneck, weight-space merging via Task Arithmetic \cite{ilharco2022editing} has emerged as a promising paradigm. Instead of storing independent dense backbones, Task Arithmetic isolates lightweight "task vectors" $\tau_k = \theta_k - \theta_{\text{base}}$ representing the parameter shift of an expert model $\theta_k$ relative to a shared pre-trained base model $\theta_{\text{base}}$. At inference, these task vectors are linearly combined and fused onto the base model: $\theta_{\text{MTL}} = \theta_{\text{base}} + \sum_k \lambda_k \tau_k$, enabling a single base network to exhibit multiple specialized capabilities simultaneously. However, even these task vectors remain fully dense (e.g., 28.7 million fine-tuned visual encoder parameters in CLIP ViT-B/32 still occupy over 115 megabytes per expert), which is still unacceptably large when deploying dozens of task-specific experts on memory-bounded edge hardware.

\begin{figure}[t]
\centering
\vskip 0.1in
\includegraphics[width=\columnwidth]{../results/pruning_resilience_curves.png}
\caption{Multi-task average accuracy across 4 datasets as a function of the task-vector parameter retention budget $p \in \{0.05, 0.10, 0.20, 1.0\}$. Norm-preserved deterministic Uniform Pruning (NP-BTVP-U) achieves outstanding results, and both AdamW and SAM experts demonstrate an extraordinary and nearly identical level of resilience under our norm-preserving rescaling.}
\label{fig:pruning_resilience_curves}
\vskip -0.1in
\end{figure}

In this paper, we approach this challenge from a deeply pragmatic perspective: \emph{How can we compress these task vectors post-hoc to a fraction of their size so they can be stored in highly compressed sparse formats like Compressed Sparse Row (CSR), while maintaining top-tier multi-task performance without adding any runtime computational latency?} 

A simple, intuitive way to compress task vectors is magnitude-based pruning, which sets the smallest absolute parameter updates to zero. In this work, we propose \textbf{Norm-Preserved Budgeted Task-Vector Pruning (NP-BTVP)}, a post-hoc weight sparsification and merging framework. We introduce norm-preserved deterministic pruning schemes: global \textbf{Uniform Pruning (NP-BTVP-U)} and layer-wise \textbf{Adaptive Saliency-Based Pruning (NP-BTVP-S)}. Crucially, we incorporate \textbf{norm-preserving rescaling} (scaling active updates by $1/p$ as a global factor, or optionally by $1/p_l$ as a layer-wise factor) as a deterministic signal-strength preservation heuristic to counteract update norm shrinkage. We explicitly clarify that while the framework is conceptually motivated by preventing norm shrinkage (hence the name "Norm-Preserved"), applying the $1/p$ scale factor to the deterministically pruned largest parameters actually amplifies the expected $L_1$ update norm mathematically (as derived in Appendix~\ref{sec:appendix_formal_analysis}). Rather than maintaining a strict identity, this reciprocal scaling acts as a highly beneficial \emph{signal-strength boost} during multi-task fusion, preventing specialized task vectors from being drowned out by the base pre-trained model (see Section 3.3 and Appendix~\ref{sec:appendix_formal_analysis}).

We evaluate our framework across standard AdamW and Sharpness-Aware Minimization (SAM) \cite{foret2021sharpness} optimizers. Our empirical results reveal surprising, counter-intuitive insights that challenge common assumptions in model merging:
First, we find that training-stage loss landscape flatness (via SAM) does not provide an additional coordinate-aligned pruning buffer under well-converged regimes compared to standard AdamW. Rather, when coupled with our norm-preserving rescaling, both standard AdamW and flatness-aware SAM experts demonstrate an extraordinary and nearly identical level of resilience to post-hoc magnitude-based coordinate-wise pruning, maintaining high classification accuracy even under extreme sparsification budgets. This separates the geometry of dense weight merging from coordinate sparsification.
Second, we show that our deterministic norm-preserved Uniform Pruning (NP-BTVP-U) is highly competitive with advanced baselines. At a tight 90\% sparsity budget ($p=0.10$), NP-BTVP-U achieves \textbf{90.32\%} average accuracy under SAM and \textbf{90.34\%} under AdamW, performing remarkably close to the dense unpruned models (\textbf{91.00\%} and \textbf{90.94\%}) and competitively with the advanced stochastic DARE-Merging baseline, while outperforming TIES-Merging (\textbf{86.51\%} and \textbf{86.65\%}) by over \textbf{3.6\%} at half the parameter budget.
Third, we identify that while layer-wise budget allocation in NP-BTVP-S is highly intuitive, it is slightly outperformed by NP-BTVP-U because its layer-wise budget heterogeneity introduces scale instability during rescaling. This underscores the pragmatist's preference for simple, globally stable solutions. We expect this Saliency Double-Bind to be particularly pronounced in broader architectures like Large Language Models (LLMs), where severe inter-layer scale mismatch and massive structural heterogeneity would further compound budget-allocation instability.

Our contributions are as follows:
\begin{itemize}
    \item We propose NP-BTVP, a training-free weight sparsification and merging framework incorporating norm-preserving rescaling, and evaluate two schemes: global NP-BTVP-U and layer-wise NP-BTVP-S.
    \item We conduct a rigorous empirical evaluation over 3 independent seeds on 4 datasets (MNIST, FashionMNIST, CIFAR-10, SVHN) using a CLIP ViT-B/32 backbone, demonstrating that deterministic global Uniform Pruning with rescaling (NP-BTVP-U) is highly effective, performing competitively with DARE and very close to the dense unpruned baseline while reducing on-disk storage by 90\%.
    \item We demonstrate that while optimizer-guided loss landscape flatness (via SAM) is crucial for stabilizing dense multi-task weight-space merging, it does not provide an additional coordinate-aligned pruning buffer under well-converged regimes, revealing a highly valuable and counter-intuitive geometric separation.
    \item We provide an in-depth, pragmatically grounded analysis of layer-wise adaptive allocation vs. global uniform pruning under rescaling, highlighting that layer-wise budget variation introduces scale instability and illustrating the pragmatic value of simple, uniform scale-preservation.
\end{itemize}